\section{A Learning Protocol for Boolean Functions}

We consider $t$ Boolean variables $p_1,\ldots,p_t$
each of which can take the value 1 or 0 to indicate
whether the associated proposition is true or
false.  There is no assumption about the independence of the variables.  Indeed, they may be functions of each other.

A \emph{vector} is an assignment to each of the $t$
variables of a value from $\{0,1,*\}$.  The symbol $*$
denotes that a variable is \emph{undetermined}.  A vector
is \emph{total} if every variable is \emph{determined} (i.e.\ is
assigned a value from $\{0,1\}$.)  For example, the
assignment $p_1 = p_3 = 1$, $p_4 = 0$ and $p_2 = {*}$ is a
vector that is not total.

A Boolean function $F$ is a mapping from the
set of $2^t$ total vectors to $\{0,1\}$.  A Boolean
function $F$ becomes a \emph{concept} if its domain
is extended to the set of all vectors as follows:
For a vector $v$, $F(v) = 1$ if and only if $F(w) = 1$
for all total vectors $w$ that agree with $v$ on
all variables for which $v$ is determined.  The
purpose of this extension is that it permits us
not to mention the variables on which $F$ does
not depend.

Given a concept $F$ we consider an \emph{arbitrary}
probability distribution $D$ over the set of all
vectors $v$ such that $F(v) = 1$.  In other words
for each $v$ such that $F(v) = 1$ it is the case
that $D(v) \geqslant 0$.  Also $\Sigma D(v) = 1$ when summation
is over this set of vectors.  There are no other
assumed restrictions on $D$, which is intended to
describe the relative frequency with which the
positive examples of $F$ occur in nature.

What are reasonable learning protocols to
consider?  First we must avoid giving the teacher
too much power, namely the power to communicate a
program instruction by instruction.  For example,
if a premeditated sequence of vectors with repetitions could be communicated then this could be used
to encode the description of the program even if
just two such vectors were used for binary notation.
Secondly we must avoid giving the teacher what is
evidently too little power.  In particular, the
protocol must provide some typical examples of

\noindent
vectors for which $F$ is true, for otherwise, if
$F$ is true for just one vector which is total,
only an exponential search or more powerful oracles
would be able to find it.

Such considerations led us to consider the
following learning protocol as a reasonable one.
It gives the learner access to the following two
routines:

\begin{enumerate}
\item[(i)] EXAMPLES:  This has no input.  It gives
as output a vector $v$ such that $F(v) = 1$.  For
each such $v$ the probability that $v$ is output
on any single call is $D(v)$.

\item[(ii)] ORACLE($\;$):  On vector $v$ as input it
outputs 1 or 0 according to whether $F(v) = 1$ or 0.
\end{enumerate}

The first of these gives randomly chosen
positive examples of the concept being learnt.
The second provides a test for whether a vector
which the deduction procedure considers to be
critical is an example of the concept or not.  In
a real system the oracle may be a human expert, a
data base of past observations, some deduction
system, or a combination of these.

Finally we observe that our particular choice
of learning protocol was strongly influenced by
the earlier decision to deal with concepts rather
than raw Boolean functions.  If we were to deal
with the latter then many other alternatives are
equally natural.  For example, given a vector $v$
instead of asking, as we do, whether all completions of it makes the function true, we could ask
whether there exist any completions that make it
true.  This suggests natural alternative semantics
for EXAMPLES and ORACLE.  In Section 7 we shall
use such more elaborate oracles.
